\title{Observation-model uncertainty and identifiability in HgCdTe band-gap extraction}

% Replace all placeholders only after the final author list is confirmed.
\author[aff1]{Author Name\corref{cor1}}
\ead{author@example.com}
\cortext[cor1]{Corresponding author. Full postal address and telephone number required.}
\affiliation[aff1]{organization={Affiliation required},
  addressline={Street address required},
  city={City required},
  postcode={Postal code required},
  country={Country required}}

\begin{abstract}
Band-gap relations for $\mathrm{Hg}_{1-x}\mathrm{Cd}_{x}\mathrm{Te}$ are often compared at the few-meV level, yet an absorption-derived edge depends on the observation definition and specimen state. We applied a fail-closed extraction ensemble to two calibrated primary-figure infrared spectroscopic ellipsometry spectra at 300~K: $x=0.226$ with thickness 15.40~$\mu$m and $x=0.310$ with thickness 4.95~$\mu$m. The ensemble included free and fixed fractional-power models, the source-panel exponent, a provenance-bound Chu Kane-region form, and fixed absorption crossings. Fitted observation-model choice produced edge spans of 6.414 and 6.830~meV. Coherent digitization-coordinate perturbations shifted every fitted-model edge by at most 0.891~meV, so digitization uncertainty did not explain the multi-meV model dependence. Fixed absorption crossings produced wider operational envelopes and changed the nominal closest material-gap comparator. The published Seiler relation was nominally closest for every fitted-model edge, but its advantage over Hansen was only 0.177--0.255~meV. Because specimen-level composition uncertainty was unreported and both spectra came from one study, that ordering was descriptive rather than an identified material-law ranking. We formulate the measured absorption edge as a latent gap plus method, carrier, vacancy, and measurement terms, then derive an audit-grade paired $2\times2\times2$ acquisition design. The results require explicit reporting of observation definition and specimen state in meV-scale HgCdTe comparisons; they do not support a universal absorption correction, selected production edge, or universal replacement for Hansen.
\end{abstract}

\begin{keyword}
HgCdTe \sep infrared absorption \sep band gap \sep optical spectroscopy \sep uncertainty \sep identifiability \sep photodetectors
\end{keyword}
